
\section{射影态射}

在 \S17.1 中，我们重新诠释了\trAffine\trMorphism：$X\to Y$ 是一个\trAffine\trMorphism，若对 $Y$ 上某个\trQC{}\trSheafOfAlgebras{} $\mathscr{B}$，存在 $Y$-概形之间的\trIsomorphism{} $X\overset{\sim}{\longleftarrow}\mathit{Spec}\,\mathscr{B}$。我们将以类似的方式来\emph{定义}\trProjective\trMorphism{}的观念。

你或许会觉得，射影性既然是如此经典的一个概念，总该存在某种表现合理的显然定义才是。但情况并非如此：射影的可能的变式定义有很多（见 [Stacks, tag 01W8]），而且它们全都有缺陷，包括我们在此给出的那个公认定义在内（见 \S17.3.4 中的种种警示）。

\textbf{17.3.1.}\ \emph{\trhDefinition{}。}\quad 一个\trMorphism{} $\pi\colon X\to Y$ 称为\textbf{\trProjective{}的}（projective），若存在一个\trIsomorphism
\[
\begin{tikzcd}[row sep=2.4em, column sep=4.6em]
X \ar[rr, "\sim"] \ar[dr, "\pi"'] && \mathit{Proj}\,\mathscr{S}_\bullet \ar[dl] \\
& Y &
\end{tikzcd}
\]
其中 $\mathscr{S}_\bullet$ 是 $Y$ 上一个满足「在次数 $1$ 处有限生成」（\mbox{假设~17.2.1}）的\trQC{}\trSheafOfAlgebras。我们称 $X$ 为一个\textbf{\trProjective{} $Y$-\trScheme}（projective $Y$-scheme），或者说 $X$ \textbf{在 $Y$ 上是\trProjective{}的}（projective over $Y$）。利用\mbox{练习~7.4.D}，这推广了射影 $A$-\trScheme{}的观念。

\textbf{17.3.A.（练习）}

（a）\emph{（射影态射的一个有用刻画）}设 $\pi\colon X\to Y$ 是一个\trMorphism。证明：$\pi$ 是射影的，当且仅当存在 $Y$ 上一个\trFiniteType{}的\trQCSheaf{} $\mathscr{S}_1$，以及一个\trClosedEmbedding{} $\iota\colon X\hookrightarrow\mathit{Proj}_Y\Sym^\bullet\mathscr{S}_1$（在 $Y$ 上，即与到 $Y$ 的诸\trMap{}相交换）。提示：\mbox{练习~17.2.H}。

（b）\emph{（射影态射的一个有用刻画，带线丛）}设 $\mathscr{L}$ 是 $X$ 上的一个\trInvertibleSheaf，而 $\pi\colon X\to Y$ 是一个\trMorphism。证明：$\pi$ 是射影的、且 $\OO(1)\cong\mathscr{L}$，当且仅当存在 $Y$ 上一个\trFiniteType{}的\trQCSheaf{} $\mathscr{S}_1$、一个\trClosedEmbedding{} $\iota\colon X\hookrightarrow\mathit{Proj}_Y\Sym^\bullet\mathscr{S}_1$（在 $Y$ 上，即与到 $Y$ 的诸\trMap{}相交换），以及一个\trIsomorphism{} $\iota^*\OO_{\mathit{Proj}_Y\Sym^\bullet\mathscr{S}_1}(1)\overset{\sim}{\longleftarrow}\mathscr{L}$。

（c）再进一步设 $Y$ 容许一个 \S16.2.5 意义下的\trAmple\trLineBundle——每当 $Y$ 是射影的、仿射的，或更一般地是拟射影的时，情形便是如此。证明：$\pi$ 是射影的，当且仅当对某个 $n$，存在一个\trClosedEmbedding{} $X\to\mathbb{P}^n_Y$（在 $Y$ 上）。这正是 [Ha1, p.~103] 中所给出的\trProjective\trMorphism{}定义。（如果你想避开带星的 \S16.2.5，那么可以假定 $Y$ 在 $\Spec A$ 上是射影的，并使用 \S16.2.1 中丰富性的定义。这样你便处理了 $Y$ 为射影的这一重要情形，但会错过其他可能有趣的情形，比如 $Y$ 为仿射或拟射影（但不\trProper）的情形。）提示：较难的方向是正向蕴含。用（a）中那个\trFiniteType{}的\trQCSheaf{} $\mathscr{S}_1$。把 $\mathscr{S}_1$ 与 $\mathscr{M}$ 的一个足够高的幂作张量，使之有限整体生成（\mbox{定理~16.2.6}，或在\trProper{}设定下的\mbox{定理~16.2.2}），从而得到一个满射
\[
\OO_Y^{\oplus(n+1)}\twoheadrightarrow\mathscr{S}_1\otimes\mathscr{M}^{\otimes\mathrm{N}}.
\]
然后用\mbox{练习~17.2.G}。

\textbf{17.3.2.}\ \textbf{射影态射的性质。}

一个态射之为射影的这一性质，显然在\trBaseChange{}下保持，因为 $\mathit{Proj}$ 构造对\trBaseChange{}表现良好（\mbox{练习~17.2.E}）。此外，\trProjective\trMorphism{}是\trProper{}的：\trProperness{}在靶上是局部的（\mbox{定理~11.4.4}(b)），而我们此前已见到射影 $A$-\trScheme{}在 $A$ 上是\trProper{}的（\mbox{定理~11.4.5}）。特别地（由\trProperness{}的定义），\trProjective\trMorphism{}是\trSeparated{}的、\trFiniteType{}的，并且是\trUniversallyClosed{}的。

\textbf{17.3.B.（重要练习：有限态射是射影的；对照\mbox{练习~8.3.1}）}\quad 按如下思路证明有限态射是射影的。设 $Z\to X$ 是\trFinite{}的，于是 $Z\cong\mathit{Spec}\,\mathscr{B}$，其中 $\mathscr{B}$ 是 $X$ 上一个\trFiniteType{}的\trQC{}\trSheafOfAlgebras。描述一个\trSheafOfGradedAlgebras{} $\mathscr{S}_\bullet$，使得 $\mathscr{S}_0\cong\OO_X$，而对 $n>0$ 有 $\mathscr{S}_n\cong\mathscr{B}$。描述一个 $X$-\trIsomorphism{} $Z\overset{\sim}{\longleftarrow}\mathit{Proj}\,\mathscr{S}_\bullet$。

特别地，\trClosedEmbedding{}是射影的。对于态射，我们有如下的蕴含序列
\[
\text{\trClosedEmbedding}\Longrightarrow\text{\trFinite}\Longrightarrow\text{\trProjective}\Longrightarrow\text{\trProper}.
\]

\textbf{17.3.3.}\quad 我们已经证明了有限态射是射影的（\mbox{练习~17.3.B}），并且具有有限的\trFiber（\mbox{练习~8.3.J}）。我们将在\mbox{定理~18.1.6} 中证明其逆，并在其后立刻陈述向\trProper\trMorphism{}的推广（尽管我们要到\mbox{定理~28.5.2} 才会证明它）。

\textbf{17.3.4.}\ \emph{关于射影态射的种种警示。}

首先，注意 $\OO(1)$——$\mathit{Proj}$ 这个概念中的一个重要组成部分——在定义中并没有被提到。（我倒宁愿它是定义的一部分，但这并不是公认的做法。）

其次，[Ha1, p.~103] 对\trProjective\trMorphism{}给出了一个不同的定义；我们遵循的是 \trGrothendieck{}的那个更一般的定义。尽管这些定义在多数场合下是同一回事（\mbox{练习~17.3.A}(c)），你在阅读文献时仍应保持警惕。例如，有限态射在 [Ha1] 的意义下并不总是射影的，而按这里的定义它们\emph{却}是射影的（\mbox{练习~17.3.B}）。

第三，射影态射在 \S8.1 的意义下是\emph{不}\trReasonable{}的！尽管\trProjective\trMorphism{}在\trBaseChange{}下保持（\S17.3.2），「是射影的」这一性质在靶上却\emph{不}是局部的（例子见 \S24.7.7）。（若再添加\trInvertibleSheaf{} $\OO(1)$ 这份额外数据，并且在诺特的情形下，它便是局部的了；见 \S17.3.8。）此外，我们将会证明，仅当最终的靶是仿射或\trNoetherian{}的时，射影态射才在复合下保持（\mbox{练习~17.3.I}）。

但是，合理性（reasonableness）的诸多推论中有许多依然成立，正如接下来几道练习所表明的那样。我们首先要证明的是：射影态射这一类态射，尽管本性不合理，却仍然满足\trCancellationTheorem~11.1.1。

\textbf{17.3.C.（练习：射影态射的消去定理）}

（a）设 $\pi\colon X\to Y$ 是一个\trClosedEmbedding，而 $\rho\colon Y\to Z$ 是射影的。证明：$\rho\circ\pi$ 是射影的。

（b）证明：一个\trMorphism{} $\rho\colon Y\to Z$ 具有射影对角，当且仅当 $\rho$ 是\trSeparated{}的。

（c）证明射影态射的消去定理：设 $\pi\colon X\to Y$ 是一个\trMorphism，$\rho\colon Y\to Z$ 是一个\trSeparated\trMorphism（亦即，由（b），它具有射影对角），而 $\rho\circ\pi$ 是射影的。证明：$\pi$ 是射影的。（可能的提示：使用\mbox{练习~11.1.A} 的提示中的那个图。）

\textbf{17.3.D.（练习）}\quad 证明：从一个射影 $k$-\trScheme{}到一个\trSeparated{} $k$-\trScheme{}的（$\Spec k$ 上的）\trMorphism{}总是射影的。（提示：前一练习。）

\textbf{17.3.E.（练习）}\quad 设 $\pi\colon X\to Y$ 与 $\pi'\colon X'\to Y'$ 是 $S$-\trScheme{}之间的两个射影态射。证明：$\pi\times_S\pi'\colon X\times_S X'\to Y\times_S Y'$ 也是射影的。提示：鉴于射影态射的不合理本性，你不能求助于\mbox{练习~8.1.A}；请改用\mbox{练习~10.6.D}。

\textbf{17.3.5.}\ \hard\ \textbf{相对情形下的\trGlobalGeneration{}与（极）\trAmpleness。}

我们把 \S16.1 的讨论推广到相对情形，为的是给自己配上相对无基点性（relatively base-point-freeness）的语言。我们不会用到这一讨论，所以在初读时你应当跳过它。但这些观念在研究文献中会反复出现。

设 $\pi\colon X\to Y$ 是一个\trQuasicompact\trQuasiSeparated\trMorphism。若 $\mathscr{F}$ 是 $X$ 上一个\trQCSheaf，我们就说 $\mathscr{F}$ 是\textbf{相对\trGloballyGenerated}（relatively globally generated）或\textbf{相对于 $\pi$ \trGloballyGenerated}（globally generated with respect to $\pi$），若拟凝聚层之间的自然映射 $\pi^*\pi_*\mathscr{F}\to\mathscr{F}$ 是满射。要保证 $\pi_*\mathscr{F}$ 是一个\trQCSheaf，拟紧性与拟分离性是必需的\footnote{原书此处作「are needed ensure」，当为「are needed to ensure」之脱漏，译文径顺化。——译者注}（\mbox{练习~14.4.C}）。但这些假设并不算很强的限制。整体生成只在拟紧的设定中才最有用，而多数人也不会为拟分离假设所烦扰。（不重要的旁注：这些假设其实可以大幅放宽。设 $\pi\colon X\to Y$ 是\emph{\trLocallyRingedSpace}之间的一个\trMorphism——不必是概形——而不带任何其他假设，且 $\mathscr{F}$ 是 $X$ 上一个\trQCSheaf，那么若 $\OO_X$-模之间的自然映射 $\pi^*\pi_*\mathscr{F}\to\mathscr{F}$ 是满射，我们就说 $\mathscr{F}$ 是相对\trGloballyGenerated{}或相对于 $\pi$ \trGloballyGenerated。）

得益于我们的这些假设，自然映射 $\pi^*\pi_*\mathscr{F}\to\mathscr{F}$ 是拟凝聚层之间的一个\trMorphism，于是「相对整体生成」这一条件在 $Y$ 上是仿射局部的。

现在设 $\mathscr{L}$ 是 $X$ 上一个\trLocallyFreeSheaf，而 $\pi\colon X\to Y$ 是一个\trMorphism。我们说 $\mathscr{L}$ 是\textbf{相对\trBasePointFree}的（relatively base-point-free）或\textbf{相对于 $\pi$ \trBasePointFree}的（base-point-free with respect to $\pi$），若它是相对\trGloballyGenerated。

\textbf{17.3.F.（练习）}\quad 设 $\mathscr{L}$ 是 $X$ 上一个相对\trBasePointFree{}的\trLineBundle，$\pi\colon X\to Y$ 是一个\trQuasicompact\trSeparated\trMorphism，而 $\pi_*\mathscr{L}$ 在 $Y$ 上是\trFiniteType{}的。（我们日后将在 \trGrothendieck{}凝聚定理~18.9.1 中证明：若 $\pi$ 是\trProper{}的且 $Y$ 是\trNoetherian{}的，则后一陈述为真。当 $\pi$ 是射影的时，这要容易得多；见\mbox{定理~18.7.1}。我们本可以现在就下功夫证明它，但那不值得费事。）描述一个典范\trMorphism{} $\psi\colon X\to\mathit{Proj}_Y\Sym^\bullet(\pi_*\mathscr{L})$。（可能的提示：这推广了「无基点的线丛给出到射影空间的映射」这一事实，所以请把那个论证加以推广；见 \S15.2。）

我们说 $\mathscr{L}$ 是\textbf{相对\trAmple}的（relatively ample）、\textbf{$\pi$-\trAmple}的（$\pi$-ample），或者\textbf{相对于 $\pi$ \trAmple}的（relatively ample with respect to $\pi$），若对 $Y$ 的每个\trAffine{}开子集 $\Spec B$，$\mathscr{L}|_{\pi^{-1}(\Spec B)}$ 在 $\pi^{-1}(\Spec B)$ 上于 $B$ 上是\trAmple{}的；或者，等价地（由 \S16.2.5），$\mathscr{L}|_{\pi^{-1}(\Spec B)}$ 在 $\pi^{-1}(\Spec B)$ 上是（绝对）\trAmple{}的。由 \S16.2.5 的讨论，若 $\mathscr{L}$ 是\trAmple{}的\footnote{原书此处作「if $\mathscr{L}$ is ample…」，未带「$\pi$-」前缀，疑当为「$\pi$-ample」之脱漏（按绝对丰富解读则非真命题），兹照录原书。——译者注}，则 $\pi$ 必然是\trQuasicompact{}的，并且（由\mbox{定理~16.2.6}）是\trSeparated{}的；若 $\pi$ 是仿射的，则一切\trInvertibleSheaf{}都是\trAmple{}的；而若 $\pi$ 是射影的，则相应的 $\OO(1)$ 是\trAmple{}的。由\mbox{练习~16.2.M}，$\mathscr{L}$ 是 $\pi$-\trAmple{}的，当且仅当 $\mathscr{L}^{\otimes n}$ 是 $\pi$-\trAmple{}的；并且若 $Z\hookrightarrow X$ 是一个\trClosedEmbedding，则 $\mathscr{L}|_Z$ 在 $Y$ 上是\trAmple{}的。

由\mbox{定理~16.2.6}(d)，我们有一个自然的\trOpenEmbedding{} $X\hookrightarrow\mathit{Proj}_Y\oplus\pi_*\mathscr{L}^{\otimes\mathrm{d}}$。（你看得出来这个映射是什么吗？另外要当心：$\oplus\pi_*\mathscr{L}^{\otimes\mathrm{d}}$ 未必是一个有限生成的\trSheafOfGradedAlgebras，所以我们这是在通常假设之一不成立的情形下使用 $\mathit{Proj}$ 构造。）

相对整体生成与相对丰富性这些观念，最有用的场合是在\trProper{}的设定中，理由是\mbox{定理~16.2.2}。

\textbf{17.3.6.}\quad \S16.1 的许多陈述可以原封不动地照搬过来。例如，我们有下面这组结果。设 $\pi\colon X\to Y$ 是\trProper{}的，$\mathscr{F}$ 与 $\mathscr{G}$ 是 $X$ 上的\trQCSheaf，而 $\mathscr{L}$ 与 $\mathscr{M}$ 是 $X$ 上的\trInvertibleSheaf。若 $\pi$ 是仿射的，则 $\mathscr{F}$ 是相对\trGloballyGenerated（容易的\mbox{练习~16.1.B}）。若 $\mathscr{F}$ 与 $\mathscr{G}$ 都是相对\trGloballyGenerated，则 $\mathscr{F}\otimes\mathscr{G}$ 亦然（容易的\mbox{练习~16.1.C}）。若 $\mathscr{L}$ 是 $\pi$-极丰富的（\mbox{定义~17.2.2}），则它是 $\pi$-\trBasePointFree{}的（容易的\mbox{练习~16.2.B}）。若 $\mathscr{L}$ 是 $\pi$-极丰富的，而 $\mathscr{M}$ 是 $\pi$-\trBasePointFree{}的（比如说，若它是 $\pi$-极丰富的），则 $\mathscr{L}\otimes\mathscr{M}$ 是 $\pi$-极丰富的（\mbox{练习~16.2.C}）。\mbox{练习~16.2.G} 可以立即推广，以证明：若
\[
\begin{tikzcd}[row sep=2.6em, column sep=4.2em]
X \ar[rr, "\pi"] \ar[dr, "\tau"'] && Y \ar[dl, "\rho"'] \\
& S &
\end{tikzcd}
\]
中的 $\pi$ 是 $S$-概形之间的一个\trFinite\trMorphism，而 $\mathscr{L}$ 是 $Y$ 上一个 $\rho$-\trAmple{}的\trInvertibleSheaf，则 $\pi^*\mathscr{L}$ 是 $\tau$-\trAmple{}的。

出于这些陈述的本性，\S16.1 的某些陈述会需要对 $Y$ 添加拟紧性假设，或者作别的修补。例如：

\textbf{17.3.7.\ \trhTheorem{}。}\ ——\ \emph{设 $\pi\colon X\to Y$ 是\trProper{}的，$\mathscr{L}$ 是 $X$ 上的一个\trInvertibleSheaf，而 $Y$ 是\trQuasicompact{}的。则以下诸条件等价。}

\emph{(a) 对某个 $\mathrm{N}>0$，$\mathscr{L}^{\otimes\mathrm{N}}$ 是 $\pi$-极丰富的。}

\emph{(a') 对一切 $n\gg 0$，$\mathscr{L}^{\otimes n}$ 是 $\pi$-极丰富的。}

\emph{(b) 对一切\trFiniteType{}的\trQCSheaf{} $\mathscr{F}$，存在某个 $n_0$，使得对 $n\ge n_0$，$\mathscr{F}\otimes\mathscr{L}^{\otimes n}$ 都是相对\trGloballyGenerated。}

\emph{(c) \trInvertibleSheaf{} $\mathscr{L}$ 是 $\pi$-\trAmple{}的。}

\textbf{17.3.G.（练习）}\quad 用\mbox{定理~16.2.2} 证明\mbox{定理~17.3.7}。（不重要的评注：\mbox{定理~16.2.2} 所给出的证明用到了诺特假设，但按那里的表述，这些假设是可以去掉的。）

做完上面的练习之后，倘若你需要去掉 $Y$ 上的拟紧性假设，\mbox{定理~17.3.7} 的陈述该如何调整，将会是清楚的。

\textbf{17.3.H.（练习：极丰富性的一个有用的等价定义）}\quad 设 $\pi\colon X\to Y$ 是一个\trProper\trMorphism，$Y$ 是局部\trNoetherian{}的（从而 $X$ 亦然，因为 $\pi$ 是\trFiniteType{}的），而 $\mathscr{L}$ 是 $X$ 上的一个\trInvertibleSheaf。假设你知道在这一情形下 $\pi_*\mathscr{L}$ 是\trFiniteType{}的。（我们日后将如\mbox{练习~17.3.F} 所述证明这一点。诺特假设正是在此处被用到的。）证明：$\mathscr{L}$ 是 $\pi$-极丰富的，当且仅当（i）$\mathscr{L}$ 是相对\trBasePointFree{}的，并且（ii）\mbox{练习~17.3.F} 的那个典范\trMorphism{} $\psi\colon X\to\mathit{Proj}_Y\Sym^\bullet(\pi_*\mathscr{L})$ 是一个\trClosedEmbedding。由此断结：若 $\pi$ 是\trProper{}的，则相对极丰富性这一观念在 $Y$ 上是仿射局部的（可以在\emph{任一份} $Y$ 的仿射覆盖上检验）。\footnote{原书此处作「any affine cover Y」，当为「any affine cover of Y」之脱漏，译文径顺化。——译者注}

\textbf{17.3.8.}\quad 作为推论，\mbox{定理~17.3.7} 蕴含：相对丰富性这一观念在 $Y$ 上是仿射局部的（若 $\pi$ 是\trProper{}的且 $Y$ 是局部\trNoetherian{}的）。

\textbf{17.3.I.（重要而富有挑战性的练习：射影态射的复合通常是射影的）}\quad 设 $\pi\colon X\to Y$ 与 $\rho\colon Y\to Z$ 是两个射影态射，而 $Z$ 是仿射的或\trNoetherian{}的。证明：$\rho\circ\pi$ 是射影的。提示：\mbox{练习~17.3.A}(b) 中给出的射影性判别法将会派上用场。（i）先处理 $Z$ 为仿射的情形。构造如下的\trCommutativeDiagram，从而找出一个在 $Z$ 上的\trClosedEmbedding{} $X\hookrightarrow\mathbb{P}\mathscr{F}^{\oplus n}$。（$\mathbb{P}\mathscr{V}$ 的定义见 \S17.2.3，其中还包括一处关于对偶的警示。）在这个图中，所有的包含都是\trClosedEmbedding，而所有的花体字母指的都是\trFiniteType{}的\trQCSheaf。
\[
\begin{tikzcd}[row sep=3.1em, column sep=1.7em]
X \ar[r, hook] \ar[drr, "\pi"'] & \mathbb{P}\mathscr{G} \ar[r, hook, "(\dagger)"] \ar[dr] & \mathbb{P}^{n-1}_Z\times_Z Y \ar[r, hook] \ar[d] & \mathbb{P}^{n-1}_Z\times_Z\mathbb{P}\mathscr{F} \ar[r, hook, "\mathrm{Segre}", "\text{参\mbox{练习~10.6.C、10.6.D}}"'] \ar[d] & \mathbb{P}(\mathscr{F}^{\oplus n}) \ar[ddll] \\
&& Y \ar[r, hook] \ar[d, "\rho"'] & \mathbb{P}\mathscr{F} \ar[dl] & \\
&& Z & &
\end{tikzcd}
\]
按如下方式构造这个\trClosedEmbedding{} $(\dagger)$。设 $\mathscr{M}$ 是 $Y$ 在 $Z$ 上的那个极丰富线丛。于是 $\mathscr{M}$ 是\trAmple{}的，从而由\mbox{定理~16.2.2}，对 $\mathrm{m}\gg 0$，$\mathscr{G}\otimes\mathscr{M}^{\otimes\mathrm{m}}$ 由有限多个\trGlobalSections{}生成。设 $\OO_Y^{\oplus n}\twoheadrightarrow\mathscr{G}\otimes\mathscr{M}^{\otimes\mathrm{m}}$ 是相应的满射。这就诱导出一个\trClosedEmbedding{} $\mathbb{P}(\mathscr{G}\otimes\mathscr{M}^{\otimes\mathrm{m}})\hookrightarrow\mathbb{P}^{n-1}_Y$。但 $\mathbb{P}(\mathscr{G}\otimes\mathscr{M}^{\otimes\mathrm{m}})\cong\mathbb{P}\mathscr{G}$（\mbox{练习~17.2.G}），而 $\mathbb{P}^{n-1}_Y=\mathbb{P}^{n-1}_Z\times_Z Y$。（ii）展开这个图以证明：（对 $Z$ 仿射的情形）若 $\mathscr{L}$ 是 $\pi$-极丰富的且 $\mathscr{M}$ 是 $\rho$-极丰富的，则对 $\mathrm{m}\gg 0$，$\mathscr{L}\otimes\pi^*\mathscr{M}^{\otimes\mathrm{m}}$ 是 $(\rho\circ\pi)$-极丰富的。然后，借助\mbox{练习~17.3.H}，通过用有限多个仿射开子集覆盖 $Z$ 来处理一般情形。

\textbf{17.3.9.}\ \extrahard\ \textbf{\trAmple\trVectorBundle。}\quad \textbf{\trAmple\trVectorBundle}（ample vector bundle）的观念在文献的某些部分中是有用的，所以我们在此给出它的定义，尽管我们不会用到这个观念。设 $\mathscr{E}$ 是\trProper{} $A$-\trScheme{} $X$ 上的一个\trLocallyFreeSheaf。若 $\OO_{\mathit{Proj}_X\Sym^\bullet\mathscr{E}/X}(1)$ 是一个\trAmple{}的\trInvertibleSheaf，则称 $\mathscr{E}$ 是\textbf{\trAmple{}的}（ample）。特别地，利用\mbox{练习~17.2.G}，你可以验证：一个\trInvertibleSheaf{}作为\trLocallyFreeSheaf（按这个定义）是\trAmple{}的，当且仅当它作为\trInvertibleSheaf（\mbox{定义~16.2.1}）是\trAmple{}的，从而免于一场记号危机。（这个\trProper{}假设是可以放宽的；之所以把它收进来，只是因为 \mbox{定义~16.2.1} 的丰富性仅对\trProper{}概形而言。）

\textbf{17.3.10.}\ \extrahard\ \textbf{拟射影态射。}\quad 与射影及拟射影 $A$-\trScheme（\S4.5.10）相类比，我们可以定义拟射影态射。\emph{若 $Y$ 是\trQuasicompact{}的}，我们就说 $\pi\colon X\to Y$ 是一个\textbf{拟射影态射}（quasiprojective morphism），若 $\pi$ 可以表达成一个到某个在 $Y$ 上射影的概形的\trQuasicompact\trOpenEmbedding。这并不是一个多么好的观念，我们也不会用到它。（拟射影态射的一般定义略显微妙——见 [Gr-EGA, II.5.3]——而我们用不着它。）

\textbf{17.3.11.}\ \extrahard\ \textbf{\trQuasiaffine{}态射。}

既然我们已经引入了拟射影态射，就不妨顺带简短地引入\trQuasiaffine{}态射（以及\trQuasiaffine{}概形），因为某些读者或许有用到它们的场合。这些观念中有许多本可以早在很久以前就引入，但由于我们永远不会用到它们，我们索性在此一并处理。

一个\trScheme{} $X$ 是\textbf{\trQuasiaffine{}的}（quasiaffine），若它容许一个到某个\trAffineScheme{}的\trQuasicompact\trOpenEmbedding。这蕴含 $X$ 是\trQuasicompact{}且\trSeparated{}的。若 $X$ 是\trNoetherian{}的（对多数人而言这是最要紧的情形），则任何\trOpenEmbedding{}当然是自动\trQuasicompact{}的。

\textbf{17.3.J.（练习）}\quad 证明：$X$ 是\trQuasiaffine{}的，当且仅当典范映射 $X\to\Spec\Gamma(X,\OO_X)$（定义见\mbox{练习~7.3.G} 及其后那段话）是一个\trQuasicompact\trOpenEmbedding。于是，\trQuasiaffine{}概形随之自带一个\emph{典范的}到\trAffineScheme{}的\trQuasicompact\trOpenEmbedding。提示：为方便计，令 $A=\Gamma(X,\OO_X)$。设 $X\hookrightarrow\Spec R$ 是一个\trQuasicompact\trOpenEmbedding。我们希望证明 $X\to\Spec A$ 是一个\trQuasicompact\trOpenEmbedding。把 $X\hookrightarrow\Spec R$ 经由 $X\to\Spec A\to\Spec R$ 来分解。对任意选定的点 $p\in X$，证明 $X\to\Spec A$ 在其某个\trOpen\trNeighborhood{}中是一个开嵌入，做法如下。选 $r\in R$ 使得 $p\in\mathrm{D}(r)\subset X$。注意：若 $X_r=\{q\in X: r(q)\neq 0\}$，则 $\Gamma(X_r,\OO_X)=\Gamma(X,\OO_X)_r$（由\mbox{练习~15.4.N}，这里用到了 $X$ 是\trQuasicompact{}且\trQuasiSeparated{}的这一事实）。用它来证明映射 $X_r\to\Spec A_r$ 是一个\trIsomorphism。

不难证明：\emph{$X$} 是\trQuasiaffine{}的，当且仅当 $\OO_X$ \emph{是}\trAmple{}的，但我们不会用到这一事实。

一个\trMorphism{} $\pi\colon X\to Y$ 称为\textbf{\trQuasiaffine{}的}，若 $Y$ 的每个\trAffine{}开子集的\trInverseImage{}都是一个\trQuasiaffine{}概形。由\mbox{练习~17.3.J}，这等价于：$\pi$ 是\trQuasicompact{}且\trSeparated{}的，并且自然映射 $X\to\mathit{Spec}\,\pi_*\OO_X$ 是一个\trQuasicompact\trOpenEmbedding。这蕴含拟仿射性（quasiaffineness）这一观念在靶上是局部的（可以在一个\trOpenCover{}上检验），并且在靶上还是仿射局部的（可以选取一个仿射覆盖，然后检验这些开集的原像都是\trQuasiaffine{}的）。\trQuasiaffine{}态射在\trBaseChange{}下保持：若 $Y$ 上的一个\trMorphism{} $X\hookrightarrow Z$ 是到某个在 $Y$ 上仿射的概形的\trQuasicompact\trOpenEmbedding，则对任意 $W\to Y$，$X\times_Y W\hookrightarrow Z\times_Y W$ 是到某个在 $W$ 上仿射的概形的\trQuasicompact\trOpenEmbedding。（有趣的是，\mbox{练习~17.3.J} \emph{并}不是用来证明这一\trBaseChange{}性质的正确工具。）

容易验证，\trQuasiaffine{}态射在复合下保持 [Stacks, tag 01SN]（于是\trQuasiaffine{}态射构成一个「\trReasonable」类）。从而，\trQuasicompact{}的\trLocallyClosedEmbedding{}都是\trQuasiaffine{}的。若 $X$ 是仿射的，则 $X\to Y$ 是\trQuasiaffine{}的，当且仅当它是\trQuasicompact{}的（因为 $Y$ 的任何仿射开子集的原像都是仿射概形 $X$ 的一个开子集）。特别地，由针对\trQuasicompact{}态射的\trCancellationTheorem~11.1.1，从仿射概形出发到拟分离概形的任何\trMorphism{}都是\trQuasiaffine{}的。
